\documentclass[aspectratio=169,11pt]{beamer}

\usetheme{AnnArbor}
\usecolortheme{dolphin}
\usefonttheme{professionalfonts}
\setbeamertemplate{navigation symbols}{}
\setbeamertemplate{caption}[numbered]
\setbeamertemplate{itemize items}[circle]
\setbeamertemplate{enumerate items}[default]

\usepackage{fontspec}
\setsansfont{DejaVu Sans}
\setmonofont{DejaVu Sans Mono}
\usepackage{amsmath,amssymb}
\usepackage{booktabs}
\usepackage{array}
\usepackage{tabularx}
\usepackage{adjustbox}
\usepackage{listings}
\usepackage{xcolor}
\usepackage{tikz}
\usetikzlibrary{arrows.meta,positioning,shapes.geometric,fit,shadows}
\usepackage{hyperref}
\hypersetup{colorlinks=true,linkcolor=dbblue,urlcolor=dbblue}

% Palette mau thuong hieu
\definecolor{dbblue}{RGB}{20,74,135}
\definecolor{dblight}{RGB}{235,243,255}
\definecolor{dbgreen}{RGB}{28,120,90}
\definecolor{dborange}{RGB}{210,120,30}
\definecolor{codebg}{RGB}{246,248,250}
\definecolor{codekw}{RGB}{0,80,180}
\definecolor{codestr}{RGB}{180,60,20}
\definecolor{codecomm}{RGB}{100,110,120}

\setbeamercolor{structure}{fg=dbblue}
\setbeamercolor{title}{fg=white,bg=dbblue}
\setbeamercolor{frametitle}{fg=dbblue}
\setbeamercolor{block title}{fg=white,bg=dbblue}
\setbeamercolor{block body}{bg=dblight,fg=black}
\setbeamercolor{block title example}{fg=white,bg=dbgreen}
\setbeamercolor{block body example}{bg=dblight!60!white,fg=black}
\setbeamercolor{block title alerted}{fg=white,bg=dborange}
\setbeamercolor{block body alerted}{bg=dblight!40!white,fg=black}

\lstset{
  basicstyle=\ttfamily\scriptsize,
  keywordstyle=\color{codekw}\bfseries,
  stringstyle=\color{codestr},
  commentstyle=\color{codecomm}\itshape,
  backgroundcolor=\color{codebg},
  frame=single,
  rulecolor=\color{dbblue!30},
  breaklines=true,
  showstringspaces=false,
  keepspaces=true,
  aboveskip=1.5pt,
  belowskip=1.5pt,
  language=SQL,
  morekeywords={CONSTRAINT,PRIMARY,KEY,FOREIGN,REFERENCES,NOT,NULL,UNIQUE,CHECK,DEFAULT,CASCADE,RESTRICT,SET,ACTION,NO,BETWEEN,ENGINE,AUTO_INCREMENT,DROP,ALTER,TABLE,ADD},
  tabsize=2
}

\newcommand{\kw}[1]{\textcolor{dbblue}{\textbf{#1}}}
\newcommand{\smallgap}{\vspace{0.2em}}

\title[Ràng buộc trong CREATE TABLE]{Bài giảng: Ràng buộc trong CREATE TABLE với MySQL}
\subtitle{PRIMARY KEY, FOREIGN KEY, NOT NULL, UNIQUE, CHECK, DEFAULT \& Toàn vẹn dữ liệu}
\author[Vũ Đức Minh]{Vũ Đức Minh}
\institute[NEU]{Trường Đại học Kinh tế Quốc dân (NEU), Hà Nội, Việt Nam}
\date{\today}

\begin{document}

%------------------------------------------------
% Slide 1: Trang bia
\begin{frame}
  \titlepage
\end{frame}

%------------------------------------------------
% Slide 2: Muc tieu bai giang
\begin{frame}{Mục tiêu bài giảng}
Sau khi hoàn thành bài học này, người học có thể:
\begin{itemize}\setlength{\itemsep}{1.5pt}
  \item Giải thích được bản chất và vai trò của \kw{Constraint} trong bảo đảm toàn vẹn dữ liệu.
  \item Khai báo thành thạo khóa chính đơn và khóa chính ghép bằng \kw{PRIMARY KEY}.
  \item Thiết lập mối quan hệ tham chiếu bảng bằng \kw{FOREIGN KEY \dots\ REFERENCES \dots}.
  \item Cấu hình chuẩn xác hành vi tham chiếu: \kw{CASCADE, RESTRICT, SET NULL, NO ACTION}.
  \item Hiểu rõ cơ chế vận hành và rủi ro khi tạm tắt \kw{FOREIGN\_KEY\_CHECKS}.
  \item Áp dụng các ràng buộc nghiệp vụ: \kw{NOT NULL}, \kw{UNIQUE}, \kw{CHECK} và \kw{DEFAULT}.
  \item Đặt tên ràng buộc theo \kw{quy ước chuẩn chuyên nghiệp} để dễ bảo trì schema.
  \item Kiểm tra định nghĩa cấu trúc bảng và constraint bằng \kw{SHOW CREATE TABLE}.
\end{itemize}
\end{frame}

%------------------------------------------------
% Slide 3: Noi dung bai hoc
\begin{frame}{Nội dung bài học}
\begin{columns}[t]
  \begin{column}{0.48\textwidth}
    \begin{block}{Phần 1: Toàn vẹn \& Khóa chính/ngoại}
      \begin{itemize}\setlength{\itemsep}{1.5pt}
        \item 1. Toàn vẹn dữ liệu \& Khái niệm Constraint
        \item 2. Quy ước đặt tên ràng buộc chuẩn
        \item 3. PRIMARY KEY: Khóa đơn \& Khóa ghép
        \item 4. FOREIGN KEY: Quan hệ cha - con
        \item 5. Hành vi tham chiếu (ON DELETE/UPDATE)
        \item 6. Kỹ thuật tắt mở FOREIGN\_KEY\_CHECKS
      \end{itemize}
    \end{block}
  \end{column}
  \begin{column}{0.48\textwidth}
    \begin{block}{Phần 2: Ràng buộc nghiệp vụ \& Vận dụng}
      \begin{itemize}\setlength{\itemsep}{1.5pt}
        \item 7. NOT NULL: Bắt buộc nhập liệu
        \item 8. UNIQUE: Chống trùng lặp dữ liệu
        \item 9. CHECK: Ràng buộc miền giá trị logic
        \item 10. DEFAULT: Giá trị mặc định tự sinh
        \item 11. Schema mẫu tổng hợp đa ràng buộc
        \item 12. Bẫy lỗi thường gặp, Quiz \& Bài tập
      \end{itemize}
    \end{block}
  \end{column}
\end{columns}
\end{frame}

%------------------------------------------------
\section{Khái niệm \& Quy ước đặt tên}
% Slide 4: Constraint la gi?
\begin{frame}{1. Constraint là gì? Ba mức độ toàn vẹn dữ liệu}
\begin{block}{Định nghĩa}
  \textbf{Constraint (Ràng buộc)} là các quy tắc được định nghĩa trên bảng hoặc cột mà Hệ quản trị CSDL (\textit{DBMS}) tự động kiểm tra nhằm bảo đảm dữ liệu luôn \kw{chính xác, hợp lệ và nhất quán}.
\end{block}

\smallgap
\begin{columns}[t]
  \begin{column}{0.32\textwidth}
    \begin{block}{Toàn vẹn thực thể}
      \begin{itemize}\setlength{\itemsep}{1.5pt}
        \item Mỗi dòng có định danh duy nhất.
        \item Do \kw{PRIMARY KEY} đảm bảo.
        \item Cấm trùng lặp và cấm NULL.
      \end{itemize}
    \end{block}
  \end{column}
  \begin{column}{0.32\textwidth}
    \begin{block}{Toàn vẹn tham chiếu}
      \begin{itemize}\setlength{\itemsep}{1.5pt}
        \item Liên kết các bảng luôn hợp lệ.
        \item Do \kw{FOREIGN KEY} đảm bảo.
        \item Chặn con trỏ tới cha không tồn tại.
      \end{itemize}
    \end{block}
  \end{column}
  \begin{column}{0.32\textwidth}
    \begin{block}{Toàn vẹn miền giá trị}
      \begin{itemize}\setlength{\itemsep}{1.5pt}
        \item Dữ liệu đúng quy tắc nghiệp vụ.
        \item Do \kw{NOT NULL}, \kw{CHECK}, \kw{UNIQUE}, \kw{DEFAULT} đảm bảo.
      \end{itemize}
    \end{block}
  \end{column}
\end{columns}
\end{frame}

%------------------------------------------------
% Slide 5: Quy uoc dat ten
\begin{frame}{2. Quy ước đặt tên Constraint chuẩn chuyên nghiệp}
\begin{alertblock}{Tại sao nên đặt tên tường minh?}
  Nếu không chỉ định tên, MySQL tự sinh tên ngẫu nhiên (ví dụ \texttt{students\_chk\_1}). Khi chạy lệnh \texttt{ALTER TABLE DROP CONSTRAINT}, việc truy vết và quản lý trong script migration sẽ rất phức tạp!
\end{alertblock}

\begin{center}
\resizebox{0.96\textwidth}{!}{
\begin{tabular}{llll}
\toprule
\textbf{Loại Constraint} & \textbf{Tiền tố} & \textbf{Cấu trúc mẫu} & \textbf{Ví dụ minh họa} \\
\midrule
\textbf{Primary Key} & \texttt{pk\_} & \texttt{pk\_<ten\_bang>} & \texttt{pk\_students}, \texttt{pk\_order\_lines} \\
\textbf{Foreign Key} & \texttt{fk\_} & \texttt{fk\_<bang\_con>\_<bang\_cha>} & \texttt{fk\_employees\_departments} \\
\textbf{Unique Key} & \texttt{uq\_} & \texttt{uq\_<bang>\_<cot>} & \texttt{uq\_users\_email}, \texttt{uq\_students\_code} \\
\textbf{Check Rule} & \texttt{chk\_} & \texttt{chk\_<bang>\_<nghiep\_vu>} & \texttt{chk\_products\_price}, \texttt{chk\_users\_age} \\
\bottomrule
\end{tabular}
}
\end{center}

\smallgap
\begin{exampleblock}{Cú pháp khai báo tường minh}
  \texttt{CONSTRAINT constraint\_name <loai\_rang\_buoc> (<danh\_sach\_cot>)}
\end{exampleblock}
\end{frame}

%------------------------------------------------
\section{PRIMARY KEY: Khóa chính}
% Slide 6: PRIMARY KEY don
\begin{frame}[fragile]{3. PRIMARY KEY đơn cột (Single-column Primary Key)}
\begin{columns}[t]
  \begin{column}{0.48\textwidth}
    \begin{block}{Đặc điểm của PRIMARY KEY}
      \begin{itemize}\setlength{\itemsep}{1.5pt}
        \item Giá trị phải \kw{duy nhất tuyệt đối} trong bảng.
        \item \kw{Không được chứa giá trị NULL}.
        \item Mỗi bảng \kw{chỉ có duy nhất một} Primary Key.
        \item Tự động tạo \kw{Clustered Index} (trong InnoDB), giúp tối ưu tốc độ tìm kiếm theo khóa.
      \end{itemize}
    \end{block}
  \end{column}
  \begin{column}{0.50\textwidth}
\begin{lstlisting}
-- Khai bao inline tren cot
CREATE TABLE departments (
    dept_id   INT PRIMARY KEY,
    dept_name VARCHAR(100) NOT NULL
);

-- Khai bao table constraint (Khuyen nghi)
CREATE TABLE departments (
    dept_id   INT,
    dept_name VARCHAR(100) NOT NULL,
    phone     VARCHAR(20),
    CONSTRAINT pk_departments 
        PRIMARY KEY (dept_id)
) ENGINE = InnoDB;
\end{lstlisting}
  \end{column}
\end{columns}
\end{frame}

%------------------------------------------------
% Slide 7: Composite PRIMARY KEY
\begin{frame}[fragile]{4. Khóa chính ghép (Composite Primary Key)}
\begin{columns}[t]
  \begin{column}{0.48\textwidth}
    \begin{block}{Bản chất của Composite Key}
      \begin{itemize}\setlength{\itemsep}{1.5pt}
        \item Gồm từ \kw{hai cột trở lên} kết hợp thành.
        \item Từng cột riêng lẻ có thể lặp lại giá trị.
        \item \kw{Tổ hợp tất cả các cột} trong khóa phải duy nhất.
        \item \textbf{Ứng dụng}: Bảng nối quan hệ nhiều - nhiều ($N-N$), ví dụ sinh viên đăng ký môn học theo kỳ.
      \end{itemize}
    \end{block}
  \end{column}
  \begin{column}{0.50\textwidth}
\begin{lstlisting}
CREATE TABLE enrollments (
    student_id  INT,
    course_id   INT,
    semester    VARCHAR(20),
    enrolled_at DATE,
    CONSTRAINT pk_enrollments 
        PRIMARY KEY (student_id, course_id, semester)
) ENGINE = InnoDB;

-- Cac dong hop le:
-- (1, 101, '2026A')
-- (1, 102, '2026A') -> student_id lap lai
-- (1, 101, '2026B') -> semester khac nhau

-- Loi Duplicate entry vi trung to hop 3 cot:
-- (1, 101, '2026A')
\end{lstlisting}
  \end{column}
\end{columns}
\end{frame}

%------------------------------------------------
\section{FOREIGN KEY: Khóa ngoại}
% Slide 8: FOREIGN KEY la gi?
\begin{frame}[fragile]{5. FOREIGN KEY: Quan hệ tham chiếu Cha - Con}
\begin{columns}[t]
  \begin{column}{0.48\textwidth}
    \begin{block}{Nguyên lý Khóa ngoại}
      \begin{itemize}\setlength{\itemsep}{1.5pt}
        \item Cột ở \kw{Bảng con} tham chiếu đến \kw{PRIMARY KEY} hoặc \kw{UNIQUE KEY} ở \kw{Bảng cha}.
        \item Chặn đứng việc tạo dữ liệu con “mồ côi” (trỏ đến bản ghi cha không tồn tại).
        \item Kiểu dữ liệu giữa cột con và cột cha phải tương thích hoàn toàn.
      \end{itemize}
    \end{block}
  \end{column}
  \begin{column}{0.50\textwidth}
\begin{lstlisting}
-- Bang cha: departments
CREATE TABLE departments (
    dept_id   INT PRIMARY KEY,
    dept_name VARCHAR(100) NOT NULL
) ENGINE = InnoDB;

-- Bang con: employees
CREATE TABLE employees (
    emp_id    INT PRIMARY KEY,
    full_name VARCHAR(100) NOT NULL,
    dept_id   INT,
    CONSTRAINT fk_employees_departments
        FOREIGN KEY (dept_id)
        REFERENCES departments(dept_id)
) ENGINE = InnoDB;
\end{lstlisting}
  \end{column}
\end{columns}
\end{frame}

%------------------------------------------------
% Slide 9: ON DELETE va ON UPDATE
\begin{frame}{6. Hành vi tham chiếu: \texttt{ON DELETE} \& \texttt{ON UPDATE}}
\begin{center}
\resizebox{0.96\textwidth}{!}{
\begin{tabular}{lp{6cm}p{6cm}}
\toprule
\textbf{Tùy chọn} & \textbf{Khi dòng ở bảng cha bị XÓA (ON DELETE)} & \textbf{Khi khóa ở bảng cha bị SỬA (ON UPDATE)} \\
\midrule
\texttt{RESTRICT} & \textbf{Chặn lệnh xóa} nếu có ít nhất 1 dòng con đang tham chiếu (mặc định an toàn) & \textbf{Chặn lệnh sửa} khóa cha nếu có dòng con đang liên kết \\
\texttt{CASCADE} & \textbf{Tự động xóa luôn} tất cả các dòng con liên quan & \textbf{Tự động cập nhật khóa mới} xuống tất cả các dòng con \\
\texttt{SET NULL} & \textbf{Đặt cột FK ở con thành NULL} (yêu cầu cột con cho phép nhận NULL) & \textbf{Đặt cột FK ở con thành NULL} \\
\bottomrule
\end{tabular}
}
\end{center}

\smallgap
\begin{alertblock}{Cảnh báo an toàn dữ liệu}
  Không lạm dụng \texttt{CASCADE} tùy tiện! Chỉ dùng \texttt{ON DELETE CASCADE} khi quan hệ là phụ thuộc sống còn (như xóa \texttt{Orders} thì xóa các dòng trong \texttt{OrderLines}). Với dữ liệu tài chính/nhân sự, luôn dùng \texttt{RESTRICT} hoặc \texttt{SET NULL}.
\end{alertblock}
\end{frame}

%------------------------------------------------
% Slide 10: ALTER TABLE quan ly FK
\begin{frame}[fragile]{7. Quản lý Khóa ngoại bằng \texttt{ALTER TABLE}}
\begin{block}{Thêm và xóa khóa ngoại trên bảng đang chạy}
  Trong quy trình phát triển thực tế, bảng thường được tạo trước, sau đó khóa ngoại mới được thêm vào thông qua các script migration.
\end{block}

\begin{lstlisting}
-- 1. Them Foreign Key moi vao bang da ton tai
ALTER TABLE employees
ADD CONSTRAINT fk_employees_departments
FOREIGN KEY (dept_id) REFERENCES departments(dept_id)
ON UPDATE CASCADE ON DELETE RESTRICT;

-- 2. Xem dinh nghia va ten constraint thuc te
SHOW CREATE TABLE employees;

-- 3. Xoa Foreign Key khi can thay doi logic nghiep vu
ALTER TABLE employees
DROP FOREIGN KEY fk_employees_departments;
\end{lstlisting}

\begin{exampleblock}{Lưu ý}
  Lệnh \texttt{DROP FOREIGN KEY} xóa ràng buộc logic nhưng InnoDB có thể vẫn giữ lại index phụ trên cột. Nếu muốn bỏ hẳn index, dùng thêm \texttt{DROP INDEX}.
\end{exampleblock}
\end{frame}

%------------------------------------------------
% Slide 11: FOREIGN_KEY_CHECKS
\begin{frame}[fragile]{8. Kỹ thuật tắt kiểm tra khóa ngoại: \texttt{FOREIGN\_KEY\_CHECKS}}
\begin{columns}[t]
  \begin{column}{0.48\textwidth}
    \begin{block}{Mục đích sử dụng}
      \begin{itemize}\setlength{\itemsep}{1.5pt}
        \item Nạp khối lượng lớn dữ liệu (\textit{bulk import}) khi bảng cha và con phụ thuộc vòng tròn.
        \item Tái tạo nhanh cấu trúc schema trong môi trường lab hoặc staging.
        \item Truncate bảng cha mà không cần gỡ bỏ quan hệ từng bảng con.
      \end{itemize}
    \end{block}
  \end{column}
  \begin{column}{0.50\textwidth}
\begin{lstlisting}
-- Kiem tra trang thai
SELECT @@FOREIGN_KEY_CHECKS;

-- Tat kiem tra khoa ngoai (Session)
SET FOREIGN_KEY_CHECKS = 0;

-- Thao tac du lieu dac biet
TRUNCATE TABLE departments;

-- Bat buoc phai bat lai ngay sau do!
SET FOREIGN_KEY_CHECKS = 1;
\end{lstlisting}
  \end{column}
\end{columns}

\begin{alertblock}{Rủi ro nghiêm trọng}
  Khi \texttt{FOREIGN\_KEY\_CHECKS = 0}, MySQL cho phép lưu dữ liệu vi phạm tham chiếu. Khi bật lại ($= 1$), MySQL \textbf{không tự rà soát sửa dữ liệu cũ}! Tuyệt đối không dùng trên production để che giấu lỗi thiết kế.
\end{alertblock}
\end{frame}

%------------------------------------------------
\section{NOT NULL \& UNIQUE}
% Slide 12: NOT NULL
\begin{frame}[fragile]{9. Ràng buộc \texttt{NOT NULL}: Bắt buộc nhập liệu}
\begin{columns}[t]
  \begin{column}{0.48\textwidth}
    \begin{block}{Nguyên lý NOT NULL}
      \begin{itemize}\setlength{\itemsep}{1.5pt}
        \item Buộc cột phải luôn có giá trị hợp lệ khi \texttt{INSERT} hoặc \texttt{UPDATE}.
        \item Tránh phát sinh giá trị \texttt{UNKNOWN} trong logic ba ngôi (\textit{three-valued logic}).
        \item Hỗ trợ tối ưu hóa lưu trữ và đánh chỉ mục.
      \end{itemize}
    \end{block}
  \end{column}
  \begin{column}{0.50\textwidth}
\begin{lstlisting}
CREATE TABLE users_not_null (
    user_id    INT PRIMARY KEY,
    full_name  VARCHAR(100) NOT NULL,
    email      VARCHAR(150) NOT NULL,
    created_at DATETIME NOT NULL
) ENGINE = InnoDB;

-- Chen hop le:
INSERT INTO users_not_null VALUES 
(1, 'Nguyen An', 'an@neu.vn', NOW());

-- Bao loi 1048 (Column cannot be null):
INSERT INTO users_not_null VALUES 
(2, NULL, 'binh@neu.vn', NOW());
\end{lstlisting}
  \end{column}
\end{columns}

\begin{exampleblock}{Phân biệt NULL và Chuỗi rỗng \texttt{''}}
  \texttt{NULL} là \kw{chưa có giá trị}. Chuỗi rỗng \texttt{''} là \kw{chuỗi có chiều dài 0}. Cột \texttt{NOT NULL} vẫn nhận chuỗi rỗng! Muốn cấm chuỗi rỗng, cần kết hợp thêm \texttt{CHECK (LENGTH(TRIM(cột)) > 0)}.
\end{exampleblock}
\end{frame}

%------------------------------------------------
% Slide 13: UNIQUE
\begin{frame}[fragile]{10. Ràng buộc \texttt{UNIQUE}: Đảm bảo tính duy nhất}
\begin{columns}[t]
  \begin{column}{0.48\textwidth}
    \begin{block}{So sánh UNIQUE vs PRIMARY KEY}
      \begin{itemize}\setlength{\itemsep}{1.5pt}
        \item \textbf{Số lượng}: Mỗi bảng chỉ có 1 PK, nhưng có thể có \kw{nhiều ràng buộc UNIQUE}.
        \item \textbf{Giá trị NULL}: Trong MySQL, cột UNIQUE cho phép chứa \kw{nhiều giá trị NULL} (trừ khi có thêm NOT NULL).
        \item Để đảm bảo vừa bắt buộc có giá trị, vừa duy nhất: dùng \kw{\texttt{NOT NULL UNIQUE}}.
      \end{itemize}
    \end{block}
  \end{column}
  \begin{column}{0.50\textwidth}
\begin{lstlisting}
CREATE TABLE room_bookings (
    booking_id   INT PRIMARY KEY,
    room_id      INT NOT NULL,
    booking_date DATE NOT NULL,
    time_slot    VARCHAR(20) NOT NULL,
    
    -- UNIQUE to hop 3 cot
    CONSTRAINT uq_room_bookings_slot
        UNIQUE (room_id, booking_date, time_slot)
) ENGINE = InnoDB;

-- Cung 1 phong trong 1 ngay tai 1 ca chi duoc dat duy nhat 1 lan!
\end{lstlisting}
  \end{column}
\end{columns}
\end{frame}

%------------------------------------------------
\section{CHECK \& DEFAULT}
% Slide 14: CHECK
\begin{frame}[fragile]{11. Ràng buộc \texttt{CHECK}: Miền giá trị logic nghiệp vụ}
\begin{columns}[t]
  \begin{column}{0.48\textwidth}
    \begin{block}{Đặc điểm của CHECK}
      \begin{itemize}\setlength{\itemsep}{1.5pt}
        \item Buộc dữ liệu phải thỏa mãn một biểu thức Boolean (trả về \texttt{TRUE} hoặc \texttt{UNKNOWN}).
        \item Hỗ trợ thực thi từ \kw{MySQL 8.0.16} trở lên.
        \item Có thể kiểm tra logic trên 1 cột hoặc \kw{so sánh tương quan giữa nhiều cột} trong dòng.
      \end{itemize}
    \end{block}
  \end{column}
  \begin{column}{0.50\textwidth}
\begin{lstlisting}
CREATE TABLE promotions (
    promo_id     INT PRIMARY KEY,
    promo_name   VARCHAR(100) NOT NULL,
    start_date   DATE NOT NULL,
    end_date     DATE NOT NULL,
    discount_pct DECIMAL(5,2) NOT NULL,
    
    -- So sanh giua 2 cot
    CONSTRAINT chk_promo_dates
        CHECK (end_date >= start_date),
        
    -- Mien gia tri phan tram
    CONSTRAINT chk_promo_discount
        CHECK (discount_pct BETWEEN 0 AND 100)
) ENGINE = InnoDB;
\end{lstlisting}
  \end{column}
\end{columns}
\end{frame}

%------------------------------------------------
% Slide 15: DEFAULT
\begin{frame}[fragile]{12. Ràng buộc \texttt{DEFAULT}: Cung cấp giá trị tự động}
\begin{columns}[t]
  \begin{column}{0.48\textwidth}
    \begin{block}{Cơ chế hoạt động của DEFAULT}
      \begin{itemize}\setlength{\itemsep}{1.5pt}
        \item Tự động điền giá trị quy định nếu lệnh \texttt{INSERT} không cung cấp cột đó.
        \item Có thể dùng hằng số (\texttt{'Pending'}, \texttt{0}) hoặc hàm thời gian (\texttt{CURRENT\_TIMESTAMP}).
        \item Nếu người dùng truyền giá trị tường minh, giá trị đó sẽ ghi đè giá trị default.
      \end{itemize}
    \end{block}
  \end{column}
  \begin{column}{0.50\textwidth}
\begin{lstlisting}
CREATE TABLE tasks (
    task_id    INT PRIMARY KEY,
    task_name  VARCHAR(150) NOT NULL,
    status     VARCHAR(20) NOT NULL DEFAULT 'Pending',
    priority   INT NOT NULL DEFAULT 3,
    created_at DATETIME NOT NULL DEFAULT CURRENT_TIMESTAMP
) ENGINE = InnoDB;

-- Chi chen task_id va task_name
INSERT INTO tasks (task_id, task_name)
VALUES (1, 'Soan de thi cuoi ky');

-- status tu dong = 'Pending'
-- priority tu dong = 3
-- created_at tu dong = thoi gian hien tai!
\end{lstlisting}
  \end{column}
\end{columns}
\end{frame}

%------------------------------------------------
% Slide 16: Tong hop toan ven
\begin{frame}[fragile]{13. Thực hành: Bản thiết kế đa ràng buộc hoàn chỉnh}
\begin{lstlisting}
CREATE TABLE course_registrations (
    student_id    INT NOT NULL,
    course_id     INT NOT NULL,
    semester      VARCHAR(20) NOT NULL DEFAULT '2026A',
    final_score   DECIMAL(4,2),
    registered_at DATETIME NOT NULL DEFAULT CURRENT_TIMESTAMP,
    
    -- 1. Composite Primary Key
    CONSTRAINT pk_course_registrations 
        PRIMARY KEY (student_id, course_id, semester),
        
    -- 2. Foreign Key tham chieu den bang sinh vien va mon hoc
    CONSTRAINT fk_cr_student FOREIGN KEY (student_id) REFERENCES students(student_id)
        ON UPDATE CASCADE ON DELETE RESTRICT,
    CONSTRAINT fk_cr_course FOREIGN KEY (course_id) REFERENCES courses(course_id)
        ON UPDATE CASCADE ON DELETE RESTRICT,
        
    -- 3. Check Constraint kiem soat diem so hop le
    CONSTRAINT chk_cr_score CHECK (final_score BETWEEN 0.00 AND 10.00)
) ENGINE = InnoDB;
\end{lstlisting}
\end{frame}

%------------------------------------------------
% Slide 17: Sai lam thuong gap
\begin{frame}{14. Các bẫy lỗi thường gặp với Constraints}
\begin{columns}[t]
  \begin{column}{0.48\textwidth}
    \begin{alertblock}{Lỗi quan hệ \& Thứ tự thực thi}
      \begin{itemize}\setlength{\itemsep}{2pt}
        \item \textbf{Tạo bảng con trước bảng cha}: Lỗi không tìm thấy cha khi tạo FK. Phải tạo cha trước, con sau.
        \item \textbf{Dùng SET NULL trên cột NOT NULL}: Xung đột logic nghiêm trọng, MySQL sẽ chặn ngay khi định nghĩa.
        \item \textbf{Lạm dụng ON DELETE CASCADE}: Vô tình xóa toàn bộ dữ liệu lịch sử hóa đơn khi người dùng bị xóa.
      \end{itemize}
    \end{alertblock}
  \end{column}
  \begin{column}{0.48\textwidth}
    \begin{alertblock}{Lỗi kiểm tra \& Cấu hình}
      \begin{itemize}\setlength{\itemsep}{2pt}
        \item \textbf{Quên bật lại FOREIGN\_KEY\_CHECKS}: Để hệ thống chạy ở trạng thái tắt kiểm tra làm dữ liệu rác tràn vào CSDL.
        \item \textbf{Nghĩ UNIQUE cấm NULL}: Cột UNIQUE không có NOT NULL sẽ cho phép vô số dòng mang giá trị NULL.
        \item \textbf{Chạy CHECK trên MySQL cũ}: Trước bản 8.0.16, MySQL chỉ nhận diện cú pháp nhưng không kiểm tra thật.
      \end{itemize}
    \end{alertblock}
  \end{column}
\end{columns}
\end{frame}

%------------------------------------------------
% Slide 18: Quiz nhanh
\begin{frame}{15. Câu hỏi ôn tập \& Củng cố}
\begin{columns}[t]
  \begin{column}{0.48\textwidth}
    \begin{block}{Câu 1: Số lượng khóa chính}
      Một bảng trong CSDL quan hệ có thể có tối đa bao nhiêu khóa chính (Primary Key)?
      \begin{itemize}
        \item A. Tùy ý bao nhiêu cũng được
        \item B. \textbf{\kw{Duy nhất 1 Primary Key}} (Chính xác!)
        \item C. Tối đa 2 (1 đơn, 1 ghép)
        \item D. Không bắt buộc giới hạn
      \end{itemize}
    \end{block}
  \end{column}
  \begin{column}{0.48\textwidth}
    \begin{block}{Câu 2: Hành vi CASCADE}
      Khi bảng cha có dòng bị xóa, tùy chọn nào tự động xóa các dòng con tương ứng?
      \begin{itemize}
        \item A. \texttt{RESTRICT}
        \item B. \texttt{SET NULL}
        \item C. \textbf{\kw{CASCADE}} (Chính xác!)
        \item D. \texttt{NO ACTION}
      \end{itemize}
    \end{block}
  \end{column}
\end{columns}
\smallgap
\begin{exampleblock}{Câu 3: Đặt tên Constraint}
  \textbf{Hỏi}: Vì sao nên tuân thủ quy ước đặt tên như \texttt{pk\_orders}, \texttt{fk\_orders\_customers}?\\
  \textbf{Đáp}: Giúp dễ dàng nhận diện loại ràng buộc, đọc log lỗi nhanh chóng và quản lý lệnh \texttt{DROP CONSTRAINT} chính xác trong các script cập nhật phiên bản.
\end{exampleblock}
\end{frame}

%------------------------------------------------
% Slide 19: Tong ket
\begin{frame}{Tổng kết bài học}
\begin{columns}[t]
  \begin{column}{0.5\textwidth}
    \begin{block}{3 Nguyên tắc cốt lõi}
      \begin{itemize}\setlength{\itemsep}{2pt}
        \item \kw{Chặn lỗi tại tầng Database}: Không bao giờ phụ thuộc 100\% vào validation ở tầng ứng dụng.
        \item \kw{Minh bạch quan hệ}: Khai báo đầy đủ Foreign Key để bảo toàn liên kết thực thể.
        \item \kw{Đầy đủ ràng buộc}: Luôn kết hợp linh hoạt \texttt{NOT NULL}, \texttt{UNIQUE}, \texttt{CHECK} và \texttt{DEFAULT}.
      \end{itemize}
    \end{block}
  \end{column}
  \begin{column}{0.46\textwidth}
    \begin{exampleblock}{Bài học tiếp theo}
      \begin{itemize}\setlength{\itemsep}{2pt}
        \item Bài tiếp theo: \textbf{Quản lý vòng đời Table trong MySQL} (\texttt{ALTER TABLE}, \texttt{RENAME}, \texttt{DROP}, \texttt{TRUNCATE} \& Generated Columns).
        \item Thực hành: Hoàn thành kịch bản lab schema \texttt{constraints\_lab}.
      \end{itemize}
    \end{exampleblock}
  \end{column}
\end{columns}
\end{frame}

\end{document}
